\section*{अध्याय \ref{ch:g8:balanced-equations} --- द्रव्यमान संरक्षण और संतुलित समीकरण}

\begin{solution}{exo:g8:balanced-equations:1}
\omterm{def:g7:chemical-reactions:reaction}{रासायनिक अभिक्रिया} में बने उत्पादों का कुल द्रव्यमान खपत हुए \omterm{def:g7:chemical-reactions:reactant}{अभिकारकों} के
कुल द्रव्यमान के बराबर होता है।
\end{solution}

\begin{solution}{exo:g8:balanced-equations:2}
गुणांक 3 है। यह 3 कार्बन \omterm{def:g7:atoms-and-molecules:atom}{परमाणुओं} और $3 \times 2 = 6$
ऑक्सीजन \omterm{def:g7:atoms-and-molecules:atom}{परमाणुओं} को दर्शाता है।
\end{solution}

\begin{solution}{exo:g8:balanced-equations:3}
हाँ: हर ओर 1 C और 2 O।
\end{solution}

\begin{solution}{exo:g8:balanced-equations:4}
\ce{2Mg + O2 -> 2MgO}: हर ओर 2 Mg और 2 O।
\end{solution}

\begin{solution}{exo:g8:balanced-equations:5}
नाइट्रोजन: बाईं ओर 2 N, इसलिए 2 \ce{NH3}; तब दाईं ओर 6 H, इसलिए 3
\ce{H2}: \ce{N2 + 3H2 -> 2NH3}।
\end{solution}

\begin{solution}{exo:g8:balanced-equations:6}
\omterm{def:g8:balanced-equations:conservation}{द्रव्यमान संरक्षण} से: $44 - 12 = \qty{32}{g}$ ऑक्सीजन।
\end{solution}

\begin{solution}{exo:g8:balanced-equations:7}
बनी हुई गैस कमरे में निकल जाती है, इसलिए तुला अब उसे नहीं तौलती। कोई
द्रव्यमान नष्ट नहीं हुआ: कम हुआ द्रव्यमान कमरे की \omterm{def:g4:air-a-mixture-of-gases:air}{वायु} में है।
\end{solution}

\begin{solution}{exo:g8:balanced-equations:8}
कार्बन: हर ओर 1। हाइड्रोजन: बाईं ओर 4, इसलिए 2 \ce{H2O}। ऑक्सीजन:
दाईं ओर $2 + 2 = 4$, इसलिए 2 \ce{O2}:
\ce{CH4 + 2O2 -> CO2 + 2H2O}।
\end{solution}

\begin{solution}{exo:g8:balanced-equations:9}
\ce{H2O} को \ce{H2O2} में बदलने से पदार्थ बदल जाता है: \ce{H2O2}
हाइड्रोजन पेरॉक्साइड है, पानी नहीं। केवल गुणांक बदले जा सकते हैं:
\ce{2H2 + O2 -> 2H2O}।
\end{solution}

\begin{solution}{exo:g8:balanced-equations:10}
$20 \times 5 = 100$ ऑक्सीजन \omterm{def:g7:atoms-and-molecules:molecule}{अणु}; $20 \times 4 = 80$ पानी के
\omterm{def:g7:atoms-and-molecules:molecule}{अणु}।
\end{solution}

\begin{solution}{exo:g8:balanced-equations:11}
2 \ce{C2H6} के साथ: 4 C, इसलिए 4 \ce{CO2}; 12 H, इसलिए 6 \ce{H2O}; दाईं
ओर ऑक्सीजन $8 + 6 = 14$, इसलिए 7 \ce{O2}:
\ce{2C2H6 + 7O2 -> 4CO2 + 6H2O}।
\end{solution}

\begin{solution}{exo:g8:balanced-equations:12}
अब भी \qty{850.0}{g}। डिब्बा बंद है: जला हुआ \qty{0.4}{g} मोम
कार्बन डाइऑक्साइड और जलवाष्प बन गया है, जो उस ऑक्सीजन के साथ, जिसे उन्होंने
लिया, अब भी डिब्बे के भीतर हैं। न कुछ अंदर आया है, न बाहर गया है।
\end{solution}

\begin{solution}{pb:g8:balanced-equations:1}
\textbf{1.} \omterm{def:g4:air-a-mixture-of-gases:air}{वायु} की ऑक्सीजन से, जो लोहे के साथ जुड़ती है:
आयरन ऑक्साइड का भार लोहे और उस ऑक्सीजन के बराबर है।

\textbf{2.} सब कुछ जार के भीतर रहता है: खपत हुई ऑक्सीजन पहले से जार में थी,
और बना हुआ ऑक्साइड वहीं रहता है। कुल द्रव्यमान नहीं
बदलता।

\textbf{3.} \omterm{def:g8:balanced-equations:conservation}{द्रव्यमान संरक्षण}: किसी अभिक्रिया में उत्पादों का कुल द्रव्यमान
\omterm{def:g7:chemical-reactions:reactant}{अभिकारकों} के कुल द्रव्यमान के बराबर होता है।

\textbf{4.} \ce{4Fe + 3O2 -> 2Fe2O3}।

\textbf{5.} \ce{3Fe + 2O2 -> Fe3O4}।

\textbf{6.} बाईं ओर: 4 Fe, $3 \times 2 = 6$ O। दाईं ओर: $2 \times 2 = 4$ Fe,
$2 \times 3 = 6$ O। संतुलित।

\textbf{7.} 3 ऑक्सीजन \omterm{def:g7:atoms-and-molecules:molecule}{अणुओं} पर 4 लोहे के \omterm{def:g7:atoms-and-molecules:atom}{परमाणु}, इसलिए $400 \times
\frac{3}{4} = 300$ ऑक्सीजन \omterm{def:g7:atoms-and-molecules:molecule}{अणु}।

\textbf{8.} $0.28 \times \frac{3}{7} = \qty{0.12}{g}$ ऑक्सीजन।

\textbf{9.} $0.28 + 0.12 = \qty{0.40}{g}$ आयरन ऑक्साइड।

\textbf{10.} हाँ: \qty{0.12}{g} चाहिए और जार में लगभग
\qty{0.5}{g} थी।

\textbf{11.} खपत हुई ऑक्सीजन जार की \omterm{def:g4:air-a-mixture-of-gases:air}{वायु} से निकल चुकी है (अब वह ठोस
ऑक्साइड के भीतर है), इसलिए जार में पहले से कम गैस है और कमरे की \omterm{def:g4:air-a-mixture-of-gases:air}{वायु} अंदर
बहती है। तब तुला का पाठ्यांक अंदर आई \omterm{def:g4:air-a-mixture-of-gases:air}{वायु} के द्रव्यमान जितना
बढ़ जाता है।

\textbf{12.} जार की \omterm{def:g4:air-a-mixture-of-gases:air}{वायु} से \qty{0.12}{g} ऑक्सीजन ली गई।
\end{solution}
